\section{Semidefinite forms and spectral spread}
\label{sec:semidefinite}

Throughout this section, $X=(X_1,\ldots,X_n)$ has independent coordinates
with Lebesgue densities $0\leq f_i\leq1$.  We write
\[
 Q=X^TAX+b^TX+c,\qquad \varphi_Q(t)=\mathbb E e^{itQ},
\]
where $A$ is real and symmetric.  The vector $b$ and scalar $c$ are
arbitrary.  Density bounds refer to the essential supremum.  Eigenvalues
of a positive semidefinite matrix are numbered in decreasing order,
including zeros.

\begin{theorem}[Semidefinite forms]\label{sd:main}
If $A$ is positive semidefinite and $\lambda_3(A)>0$, then $Q$ has a
density satisfying
\begin{equation}\label{sd:main-bound}
 \|p_Q\|_\infty\leq \frac{17100}{\lambda_3(A)}.
\end{equation}
The same conclusion holds for negative semidefinite $A$, with
$\lambda_3(-A)$ in the denominator.  The bound is independent of dimension
and rank, and allows a linear term outside the range of $A$.
\end{theorem}

\begin{remark}[Density at the lower endpoint]\label{sd:endpoint-zero}
For the homogeneous form $Q=X^TAX$ with $A\succeq0$ and
$\lambda_3(A)>0$, only three positive eigendirections are needed to
show that the density vanishes \emph{at zero}.  Let $u_1,u_2,u_3$
be orthonormal eigenvectors with positive eigenvalues
$\lambda_1,\lambda_2,\lambda_3$, and put
$Y_j=\langle u_j,X\rangle$.  Lemma~\ref{four:projection} bounds the
joint density of $(Y_1,Y_2,Y_3)$ by $2^{3/2}$.  Positivity of the
discarded spectral terms gives
\[
 Q\geq\sum_{j=1}^3\lambda_jY_j^2,
 \qquad
 \mathbb P(Q\leq t)
 \leq \frac{8\sqrt2\,\pi}{3\sqrt{\lambda_1\lambda_2\lambda_3}}
       t^{3/2}\quad(t\geq0).
\]
Thus the distribution function has right derivative zero at zero;
the continuous density furnished by
Theorem~\ref{fg:fixed-matrix-continuity} satisfies $p_Q(0)=0$.
The comparison concerns the lower tail of a nonnegative homogeneous
form.  It does not identify the law of $Q$ after spectral truncation
or prove a density bound at positive levels, and it does not apply at
zero when an affine term moves zero away from the lower endpoint.
\end{remark}

\subsection{Concentration and conditioning}
For a real random variable $Z$, define
\[
 L_Z(u)=\sup_s\mathbb P(|Z-s|\leq u),\qquad
 \alpha_1(Z)=\sup_{u>0}\frac{L_Z(u)}{\sqrt u},\qquad
 \alpha_2(Z)=\sup_{0<u\leq1}
       \frac{L_Z(u)}{u(1+\log(1/u))}.
\]

\begin{lemma}[Fourier concentration]\label{sd:fejer}
For every $u>0$,
\begin{equation}\label{sd:fejer-bound}
 L_Z(u)\leq\frac{u}{2\sin^2 1}
              \int_{|t|\leq2/u}|\varphi_Z(t)|\,dt.
\end{equation}
If $|\varphi_Z(t)|\leq B|t|^{-1/2}$ for $t\ne0$, then
$\alpha_1(Z)\leq2\sqrt2 B/\sin^2 1$.
If $|\varphi_Z(t)|\leq\min(1,B/|t|)$ and $B\geq2$, then
$\alpha_2(Z)\leq B/\sin^2 1$.
\end{lemma}

\begin{proof}
The function $K(x)=(\sin x/x)^2$, with $K(0)=1$, is at least
$\sin^2 1$ on $[-1,1]$.  Its Fourier transform, with kernel $e^{-itx}$,
is $\pi(1-|t|/2)_+$.  Fourier inversion and Fubini's theorem therefore give
\[
 \mathbb P(|Z-s|\leq u)
 \leq\frac{\mathbb E K((Z-s)/u)}{\sin^2 1}
 \leq\frac{u}{2\sin^2 1}
            \int_{|t|\leq2/u}|\varphi_Z(t)|\,dt.
\]
Integrating $B|t|^{-1/2}$ proves the first consequence.
For $0<u\leq\min(1,2/B)$, integration of $\min(1,B/|t|)$ gives
\[
 L_Z(u)\leq\frac{Bu}{\sin^2 1}
       \left(1+\log\frac{2}{Bu}\right)
 \leq\frac{Bu}{\sin^2 1}(1+\log(1/u)).
\]
For $2/B<u\leq1$, the last inequality follows from $L_Z(u)\leq1$
and $Bu>2$.
\end{proof}

\begin{lemma}[Square-root profiles]\label{sd:composition}
Suppose that, for every nonnegative Borel function $h$,
\begin{equation}\label{sd:profile}
 \mathbb E h(Q)\leq
 K\,\mathbb E\int_{\mathbb R}h(y)|y-\xi|^{-1/2}\,dy.
\end{equation}
Then $Q$ has a density.  Moreover:
\begin{enumerate}
\item $\alpha_1(Q)\leq4K$;
\item if $\alpha_1(\xi)\leq K_1$, then
$\alpha_2(Q)\leq4\sqrt2 K\max(K_1,1)$;
\item if $\alpha_2(\xi)\leq K_2$ and $K_2\geq1$, then
\begin{equation}\label{sd:composition-density}
 \|p_Q\|_\infty\leq3K\sqrt{K_2(3+\log K_2)}.
\end{equation}
\end{enumerate}
\end{lemma}

\begin{proof}
Tonelli's theorem shows that the right side of \eqref{sd:profile}
vanishes when $h$ is the indicator of a Lebesgue null set.  Thus $Q$
has a density, with
$p_Q(y)\leq K\mathbb E|y-\xi|^{-1/2}$ almost everywhere.
The integral of $|y-\xi|^{-1/2}$ over an interval of length $2u$
is at most $4\sqrt u$, proving the first assertion.

For any $d\in\mathbb R$,
\begin{equation}\label{sd:interval-kernel}
 \int_{-u}^u|y+d|^{-1/2}\,dy
 \leq4\sqrt2\,u\max(|d|,u)^{-1/2}.
\end{equation}
For $|d|\leq2u$ this follows from the preceding integral bound.
For $|d|>2u$, use $|y+d|\geq|d|/2$.
When $0<u\leq1$, layer cake gives
\begin{align*}
 \mathbb E\max(|s-\xi|,u)^{-1/2}
 &=\int_0^{u^{-1/2}}\mathbb P(|s-\xi|<t^{-2})\,dt\\
 &\leq1+K_1\int_1^{u^{-1/2}}\frac{dt}{t}\\
 &\leq\max(K_1,1)(1+\log(1/u)).
\end{align*}
Combine this with \eqref{sd:profile} and \eqref{sd:interval-kernel}.

For the last assertion, put $V=\sqrt{K_2(3+\log K_2)}\geq1$.
For every $s$,
\begin{align*}
 \mathbb E|s-\xi|^{-1/2}
 &\leq V+K_2\int_V^\infty t^{-2}(1+2\log t)\,dt\\
 &=V+K_2(3+2\log V)/V.
\end{align*}
Writing $a=3+\log K_2\geq3$ gives
$3+2\log V=a+\log a\leq2a$.
The last display is at most $3V$, proving the claim.
\end{proof}

\begin{lemma}[Positive pivots]\label{sd:pivot}
Suppose $A$ is positive semidefinite and $a=a_{ii}>0$.
Completing the square gives
$Q=a(X_i-m(X_{-i}))^2+R(X_{-i})$, with $m$ affine and $R$ quadratic.
The matrix of $R$ is the positive semidefinite Schur complement
\[
 S=A_{-i,-i}-a^{-1}A_{-i,i}A_{i,-i}.
\]
The profile \eqref{sd:profile} holds with $K=a^{-1/2}$ and $\xi=R$.
For $k\geq2$,
\begin{equation}\label{sd:schur-eigenvalue}
 \lambda_{k-1}(S)\geq\lambda_k(A).
\end{equation}
If $A$ has rank $r$, then $S$ has rank $r-1$.
\end{lemma}

\begin{proof}
For fixed $X_{-i}$, the density of $a(X_i-m)^2$ at $w>0$ is
\[
 \frac{f_i(m+\sqrt{w/a})+f_i(m-\sqrt{w/a})}{2\sqrt{aw}}
 \leq(aw)^{-1/2}.
\]
Integration and averaging prove the profile.  No independence between
$m$ and $R$ is needed.  Writing
$A=\left(\begin{smallmatrix}a&v^T\\v&D\end{smallmatrix}\right)$, the identity
\[
 (x,y)^TA(x,y)=a(x+a^{-1}v^Ty)^2+y^TSy
\]
proves positivity of $S$ and the rank assertion.
A $k$-dimensional subspace where the form of $A$ is at least
$\lambda_k(A)$ times the squared norm intersects $ax+v^Ty=0$
in dimension at least $k-1$.  Projection onto $y$ is injective on
this intersection, and on its image
$y^TSy\geq\lambda_k(A)(x^2+|y|^2)\geq\lambda_k(A)|y|^2$.
The variational principle proves \eqref{sd:schur-eigenvalue}.
\end{proof}

\subsection{The decoupling normalization}
For a $p\times q$ matrix $M$ of rank at least $k$, set
\[
 Z_k=e_k(M^TM)
 =\sum_{\substack{S\subseteq[p],\,T\subseteq[q]\\|S|=|T|=k}}
       \det(M_{S,T})^2,\qquad
 w_{S,T}=\det(M_{S,T})^2/Z_k.
\]
Let $\alpha_i=\sum_{S\ni i,T}w_{S,T}$ and
$\beta_j=\sum_{S,T\ni j}w_{S,T}$ be the row and column inclusion
probabilities.  If $M=A_{I,J}$ for disjoint coordinate sets $I,J$
and $\alpha_i\leq1/2$ for every $i\in I$, Lemma~\ref{four:volume} gives
\begin{equation}\label{sd:volume-decay}
 |\varphi_Q(t)|^2\leq
 e^{k/2}(2\pi)^k|t|^{-k}Z_k^{-1/2}.
\end{equation}
Arbitrary quadratic terms inside $I$ and $J$, and arbitrary linear
terms, are allowed.  For $k=3$, Fourier inversion gives
\begin{equation}\label{sd:volume-density}
 \|p_Q\|_\infty\leq6\sqrt e\,Z_3^{-1/6}.
\end{equation}
Indeed, with $B=(2\pi)^{3/2}e^{3/4}Z_3^{-1/4}$,
\[
 (2\pi)^{-1}\int_{\mathbb R}\min(1,B|t|^{-3/2})\,dt
 =3B^{2/3}/\pi=6\sqrt e\,Z_3^{-1/6}.
\]
The analytic input in \eqref{sd:volume-decay} is the integrated
Fourier estimate of Lemma~\ref{four:lp} and the Gaussian
Brascamp--Lieb theorem \citep{BL1976,Lieb1990}.  No pointwise decay
assumption on an individual characteristic function is used.

\subsection{A partition in Gram coordinates}
For positive semidefinite $A$ of rank $r$, choose
\[
 A=BB^T,\qquad B^TB=\Lambda
 =\operatorname{diag}(\lambda_1,\ldots,\lambda_r)>0.
\]
Write $b_i^T$ for row $i$ of $B$, so that $|b_i|^2=a_{ii}$.
For independent equiprobable signs $\varepsilon_i$, define
\[
 I=\{i:\varepsilon_i=1\},\quad J=\{i:\varepsilon_i=-1\},\quad
 N=\tfrac12\sum_i\varepsilon_i b_ib_i^T.
\]
Then $H_I=B_I^TB_I=\Lambda/2+N$ and
$H_J=B_J^TB_J=\Lambda/2-N$.

\begin{lemma}[Partition identity]\label{sd:partition-identity}
For every $g\in\mathbb R^r$,
\begin{align}
 \langle B_Ig,A_{I,J}B_Jg\rangle
 &=|\Lambda g|^2/4-|Ng|^2,\label{sd:gram-identity}\\
 |B_Ig|\,|B_Jg|&\leq g^T\Lambda g/2.\label{sd:gram-product}
\end{align}
If $K\subseteq[r]$ has cardinality $k$, $\lambda_m\geq\theta>0$
for $m\in K$, and
$T_K=\sum_{m\in K}|Ne_m|^2/\lambda_m^2\leq1/8$, then
$s_k(A_{I,J})\geq\theta/4$.
If $\max_i a_{ii}\leq d$, then
\[
 \mathbb ET_K\leq kd/(4\theta),\qquad
 \mathbb P(T_K>1/8)\leq2kd/\theta.
\]
\end{lemma}

\begin{proof}
The scalar product is $g^TH_IH_Jg$.  Expanding $H_IH_J$ gives
$\Lambda^2/4-N^2+(N\Lambda-\Lambda N)/2$, whose last term is
skew symmetric.  This proves \eqref{sd:gram-identity}.
Equation~\eqref{sd:gram-product} follows from $H_I+H_J=\Lambda$.
For $g$ in the span $G$ of $\{e_m:m\in K\}$,
Cauchy--Schwarz gives $|Ng|\leq|\Lambda g|T_K^{1/2}$.
Thus the scalar product is at least $|\Lambda g|^2/8$.
It is positive for $g\ne0$, so $B_J$ is injective on $G$.
For $x=B_Jg\ne0$,
\[
 \frac{|A_{I,J}x|}{|x|}
 \geq\frac{\langle B_Ig,A_{I,J}B_Jg\rangle}{|B_Ig|\,|B_Jg|}
 \geq\frac{|\Lambda g|^2}{4g^T\Lambda g}\geq\theta/4.
\]
The singular-value variational principle on the $k$-dimensional
subspace $B_JG$ proves the assertion.  Finally,
\[
 \mathbb E|Ne_m|^2
 =\tfrac14\sum_i b_i(m)^2|b_i|^2
 \leq\tfrac d4\sum_i b_i(m)^2=d\lambda_m/4.
\]
Sum over $m$ and apply Markov's inequality.
\end{proof}

\subsection{Volume marginals and ridge leverage}

\begin{lemma}[Volume marginals]\label{sd:ridge}
Let $G=MM^T$ have eigenvalues $\sigma_1\geq\sigma_2\geq\cdots\geq0$
and an orthonormal eigenbasis $(u_m)$.  If $\operatorname{rank}M\geq k$,
then
\begin{equation}\label{sd:marginal-spectrum}
 \alpha_i=\sum_m p_m u_m(i)^2,\qquad
 p_m=\frac{\sigma_m e_{k-1}(\sigma_{-m})}{e_k(\sigma)},
\end{equation}
where $e_0=1$ and $\sigma_{-m}$ is the list with entry $m$ removed.
Moreover $p_m\leq\min(1,k\sigma_m/\sigma_k)$.
If $0<\kappa\leq\sigma_k/k$ and $m_i=M^Te_i$, then
\begin{equation}\label{sd:ridge-bound}
 \alpha_i\leq2m_i^T(M^TM+\kappa I)^{-1}m_i.
\end{equation}
\end{lemma}

\begin{proof}
By Cauchy--Binet, the numerator of $\alpha_i$ is
$\sum_{S\ni i,\,|S|=k}\det G_{S,S}$.
Differentiate $\det(I+tD(z)G)$ with respect to $z_i$ at
$z=(1,\ldots,1)$, where $D(z)=\operatorname{diag}(z)$, and take
the coefficient of $t^k$.  The principal-minor expansion gives
this numerator, whereas spectral decomposition gives the
coefficient in
\[
 \sum_m u_m(i)^2\,t\sigma_m
                     \prod_{\ell\ne m}(1+t\sigma_\ell).
\]
This proves \eqref{sd:marginal-spectrum}.
The identity
$e_k(\sigma)=\sigma_me_{k-1}(\sigma_{-m})+e_k(\sigma_{-m})$
gives $p_m\leq1$.  For $m\leq k$ it also gives
$p_m\leq k\sigma_m/\sigma_k$.  Otherwise the list
$\mu=\sigma_{-m}$ contains $k$ entries at least $\sigma_k$, and
\[
 k e_k(\mu)
 =\sum_{|S|=k-1}\left(\prod_{j\in S}\mu_j\right)
                       \sum_{j\notin S}\mu_j
 \geq\sigma_k e_{k-1}(\mu).
\]
Indeed, each $(k-1)$-element set omits one of those $k$ entries.
This proves the remaining bound for $p_m$.
The inequality $\min(1,x)\leq2x/(1+x)$ now yields
$p_m\leq2\sigma_m/(\sigma_m+\kappa)$.
The spectral decomposition of
$M(M^TM+\kappa I)^{-1}M^T$ proves \eqref{sd:ridge-bound}.
\end{proof}

\begin{lemma}[Gram comparison]\label{sd:ridge-comparison}
For $M=A_{I,J}=B_IB_J^T$, $i\in I$, and $\kappa>0$,
\begin{equation}\label{sd:ridge-gram}
 m_i^T(M^TM+\kappa I)^{-1}m_i
 \leq b_i^T(H_I+\kappa\Lambda^{-1})^{-1}b_i .
\end{equation}
\end{lemma}

\begin{proof}
The identity
$(B_JH_IB_J^T+\kappa I)B_J=B_J(H_IH_J+\kappa I)$ gives
\[
 m_i^T(M^TM+\kappa I)^{-1}m_i
 =b_i^TH_J(H_IH_J+\kappa I)^{-1}b_i .
\]
If $H_J>0$, the middle matrix is
$(H_I+\kappa H_J^{-1})^{-1}$.
Since $H_J\leq\Lambda$, order reversal under inversion gives
\eqref{sd:ridge-gram}.  If $H_J$ is singular, replace it in this
matrix identity by $H_J+\varepsilon I\leq\Lambda+\varepsilon I$
and let $\varepsilon\downarrow0$.
Continuity at a positive semidefinite matrix $T$ follows from the
expression
$T^{1/2}(T^{1/2}H_IT^{1/2}+\kappa I)^{-1}T^{1/2}$.
\end{proof}

\begin{proposition}[A semidefinite partition]\label{sd:psd-partition}
For $k\in\{1,2,3\}$, let $\theta=\lambda_k(A)>0$, with $A$
positive semidefinite, and put
$\gamma_1=128$, $\gamma_2=300$, $\gamma_3=640$.
If $\max_i a_{ii}\leq\theta/\gamma_k$, some partition
$[n]=I\sqcup J$ satisfies
\[
 s_k(A_{I,J})\geq\theta/4,\qquad
 \max_{i\in I}\alpha_i
 \leq(64k/\theta)\max_i a_{ii}\leq1/2.
\]
Uniformly in the linear and constant terms, this implies
\begin{align}
 k=1:\quad&
 |\varphi_Q(t)|\leq2e^{1/4}\sqrt{2\pi}\,
                         (\theta|t|)^{-1/2},\label{sd:psd-k1}\\
 k=2:\quad&
 |\varphi_Q(t)|\leq8\pi\sqrt e\,(\theta|t|)^{-1},
                         \label{sd:psd-k2}\\
 k=3:\quad&
 \|p_Q\|_\infty\leq24\sqrt e/\theta.\label{sd:psd-k3}
\end{align}
\end{proposition}

\begin{proof}
Put $d=\max_i a_{ii}$.  Split the Gram-coordinate space into its
first $k$ coordinates, denoted by the subscript $b$, and the
remaining coordinates, denoted by $s$.  These subscripts denote
blocks, not the linear coefficient of $Q$.  Set
\[
 \kappa=\theta^2/(16k),\qquad \rho=\kappa/\theta=\theta/(16k),
 \qquad \ell_i=|\Lambda_b^{-1/2}b_i^b|^2.
\]
Then $\sum_i\ell_i=k$ and $\ell_i\leq d/\theta$.
Consider the events
\begin{align*}
 \mathcal S:\quad&
 \sum_{m=1}^k |Ne_m|^2/\lambda_m^2\leq1/8,\\
 \mathcal E_1:\quad&
 \|\Lambda_b^{-1/2}N_{bb}\Lambda_b^{-1/2}\|_F^2\leq1/64,\\
 \mathcal E_2:\quad&
 (2/\rho)\|N_{sb}\Lambda_b^{-1/2}\|_F^2\leq1/8.
\end{align*}
An empty $s$ block contributes zero.
On $\mathcal S$, Lemma~\ref{sd:partition-identity} gives
$s_k(M)\geq\theta/4$, so $\kappa\leq s_k(M)^2/k$.
Lemmas~\ref{sd:ridge} and \ref{sd:ridge-comparison} give
$\alpha_i\leq2b_i^T(H_I+\kappa\Lambda^{-1})^{-1}b_i$.

For $T>0$, completing a square gives the block inequality
\begin{equation}\label{sd:block-comparison}
 \begin{pmatrix}P&R\\R^T&T\end{pmatrix}
 \geq
 \begin{pmatrix}P-2RT^{-1}R^T&0\\0&T/2\end{pmatrix}.
\end{equation}
The quadratic form of the difference is
$2|T^{-1/2}R^Tx+T^{1/2}y/2|^2$.
As $\Lambda_s\leq\theta I$ and $(H_I)_{ss}\geq0$,
\eqref{sd:block-comparison}, with $T=\rho I$, yields
\[
 H_I+\kappa\Lambda^{-1}
 \geq
 \left(\Lambda_b/2+N_{bb}-(2/\rho)N_{bs}N_{sb}\right)
       \oplus(\rho/2)I_s.
\]
On $\mathcal E_1$, $N_{bb}\geq-\Lambda_b/8$.
On $\mathcal E_2$, conjugating by $\Lambda_b^{-1/2}$ and bounding
operator norm by trace gives
$(2/\rho)N_{bs}N_{sb}\leq\Lambda_b/8$.
Therefore
$H_I+\kappa\Lambda^{-1}\geq\Lambda_b/4\oplus(\rho/2)I_s$ and
\[
 \alpha_i\leq8\ell_i+(4/\rho)|b_i^s|^2
 \leq(64k/\theta)|b_i|^2=(64k/\theta)a_{ii}.
\]

It remains to show that these events can hold simultaneously.
Lemma~\ref{sd:partition-identity} gives
$\mathbb P(\mathcal S^c)\leq2kd/\theta$.
Independence and centering of the signs also give
\begin{align*}
 \mathbb E\|\Lambda_b^{-1/2}N_{bb}\Lambda_b^{-1/2}\|_F^2
 &=\tfrac14\sum_i\ell_i^2\leq kd/(4\theta),\\
 \mathbb E\|N_{sb}\Lambda_b^{-1/2}\|_F^2
 &=\tfrac14\sum_i|b_i^s|^2\ell_i\leq kd/4.
\end{align*}
Markov's inequality gives failure bounds $16kd/\theta$ and
$64k^2d/\theta$ for $\mathcal E_1$ and $\mathcal E_2$.
The total is at most
\[
 (18k+64k^2)/\gamma_k
 =
 \begin{cases}82/128,&k=1,\\292/300,&k=2,\\630/640,&k=3,\end{cases}
\]
which is strictly below one.  The marginal bound is at most
$1/2$ for each stated $\gamma_k$.
Since $Z_k\geq\prod_{j=1}^k s_j(M)^2\geq(\theta/4)^{2k}$,
\eqref{sd:volume-decay}--\eqref{sd:volume-density} prove
\eqref{sd:psd-k1}--\eqref{sd:psd-k3}.
\end{proof}

\subsection{Completion of the semidefinite proof}

\begin{proposition}[The first two levels]\label{sd:psd-levels}
For positive semidefinite $A$ and arbitrary $b,c$:
\begin{enumerate}
\item if $\lambda_1(A)\geq1$, then
$\alpha_1(Q)\leq K_1:=4\sqrt{128}<45.26$;
\item if $\lambda_2(A)\geq1$, then
$\alpha_2(Q)\leq K_2:=4\sqrt{600}\,K_1<4435$.
\end{enumerate}
\end{proposition}

\begin{proof}
If some $a_{ii}\geq1/128$, Lemmas~\ref{sd:pivot} and
\ref{sd:composition}(1) give the first bound.
Otherwise Proposition~\ref{sd:psd-partition}, with $k=1$,
and Lemma~\ref{sd:fejer} give
\[
 \alpha_1(Q)\leq
 (2\sqrt2/\sin^2 1)\,2e^{1/4}\sqrt{2\pi}<25.72<K_1.
\]
For the second assertion, a coefficient $a_{ii}\geq1/300$ leaves
a positive semidefinite Schur complement with an eigenvalue at least
one.  Apply the first assertion to its quadratic form, which is the
random center in the profile.  Lemma~\ref{sd:composition}(2)
gives $\alpha_2(Q)\leq4\sqrt2\sqrt{300}\,K_1=K_2$.
If there is no such coefficient,
\eqref{sd:psd-k2} and Lemma~\ref{sd:fejer} give
$\alpha_2(Q)\leq8\pi\sqrt e/\sin^2 1<58.53<K_2$.
\end{proof}

\begin{proof}[Proof of Theorem~\ref{sd:main}]
Suppose first that $\lambda_3(A)\geq1$.  A coefficient
$a_{ii}\geq1/640$ leaves a positive semidefinite Schur complement
with at least two eigenvalues at least one.
Apply Proposition~\ref{sd:psd-levels}(2) to its quadratic form and
Lemma~\ref{sd:composition}(3) to obtain
\[
 \|p_Q\|_\infty\leq
 3\sqrt{640}\sqrt{K_2(3+\log K_2)}<17100.
\]
Here $K_2=4\sqrt{600}\,4\sqrt{128}$ is used without rounding.
If every diagonal coefficient is below $1/640$,
Proposition~\ref{sd:psd-partition}, with $k=3$, instead gives
$\|p_Q\|_\infty\leq24\sqrt e$.
For $\lambda_3(A)=\theta>0$, apply this argument to $Q/\theta$
and use density scaling.  Applying the result to $-Q$ proves
the negative semidefinite case.
\end{proof}

\subsection{Indefinite forms with a small absolute spectral diagonal}
For a real symmetric $A$, write $|A|=(A^2)^{1/2}$ and let
$|\lambda|_{(1)}(A)\geq|\lambda|_{(2)}(A)\geq\cdots$
be its eigenvalue moduli.  In particular $(|A|)_{ii}\geq|a_{ii}|$.

\begin{theorem}[Absolute spectral spread]\label{sd:indefinite}
Let $k\in\{1,2,3\}$, set $\theta=|\lambda|_{(k)}(A)>0$, and put
$\gamma'_1=130$, $\gamma'_2=360$, $\gamma'_3=680$.
If
\begin{equation}\label{sd:absolute-spread}
 \max_i(|A|)_{ii}\leq\theta/\gamma'_k,
\end{equation}
some partition $[n]=I\sqcup J$ satisfies
$s_k(A_{I,J})\geq23\theta/80$ and
$\max_{i\in I}\alpha_i\leq1/2$.
Consequently, for every $b,c$,
\begin{equation}\label{sd:indefinite-fourier}
 |\varphi_Q(t)|\leq
 e^{k/4}(2\pi)^{k/2}(23\theta|t|/80)^{-k/2}.
\end{equation}
For $k=3$, $Q$ has a density and
\begin{equation}\label{sd:indefinite-density}
 \|p_Q\|_\infty\leq \frac{480\sqrt e}{23\theta}
 <\frac{35}{\theta}.
\end{equation}
\end{theorem}

\begin{proof}
On the range of $A$, factor
\[
 A=B\Sigma B^T,\qquad B^TB=\Lambda>0,\qquad
 \Sigma=\operatorname{diag}(\pm1),
\]
where the entries of $\Lambda$ are the nonzero eigenvalue moduli
in decreasing order.  Thus $BB^T=|A|$ and
$|b_i|^2=(|A|)_{ii}\leq d:=\max_i(|A|)_{ii}$.
Define the random partition and $N$ as before.  Then
$H_I=\Lambda/2+N$, $H_J=\Lambda/2-N$, and
$M=A_{I,J}=B_I\Sigma B_J^T$.
Let $G$ be the first $k$ Gram coordinates.
Decompose $g\in G$ as $g=g_++g_-$ according to the signs in $\Sigma$.
Direct multiplication gives
\begin{align}
 \langle B_I\Sigma g,MB_Jg\rangle
 &=|\Lambda g|^2/4-g_-^T\Lambda Ng_+
      -g_+^T\Lambda Ng_-
      -(\Sigma N\Sigma g)^TNg,\label{sd:signed-identity}\\
 |B_I\Sigma g|^2+|B_Jg|^2
 &=g^T\Lambda g-4g_+^TNg_-.\label{sd:signed-norm}
\end{align}
Suppose
\begin{equation}\label{sd:signed-events}
 T:=\sum_{m=1}^k |Ne_m|^2/\lambda_m^2\leq\eta^2,\qquad
 \|\Lambda_b^{-1/2}N_{bb}\Lambda_b^{-1/2}\|_F\leq\eta',
\end{equation}
where $b$ denotes the first $k$ Gram coordinates.
Then $|Ng|\leq\eta|\Lambda g|$ for $g\in G$.
The sum of the two middle terms in \eqref{sd:signed-identity}
has modulus at most $\eta|\Lambda g|^2$, by applying this bound
separately to $g_+$ and $g_-$ and then using $2ab\leq a^2+b^2$.
The last term has modulus at most $2\eta^2|\Lambda g|^2$
(the smaller coefficient $\eta^2$ would also suffice).
Since
$|g_+^TNg_-|\leq
\eta'|\Lambda^{1/2}g_+|\,|\Lambda^{1/2}g_-|$,
\eqref{sd:signed-norm} gives
\[
 |B_I\Sigma g|\,|B_Jg|
 \leq(1/2+\eta')g^T\Lambda g.
\]
If $1/4-\eta-2\eta^2>0$, positivity of the scalar product makes
$B_J$ injective on $G$.  Testing against $B_I\Sigma g$ and applying
the singular-value variational principle gives
\begin{equation}\label{sd:signed-sv}
 s_k(M)\geq
 \theta\,\frac{1/4-\eta-2\eta^2}{1/2+\eta'}.
\end{equation}
Take $\eta=1/16$ and $\eta'=1/8$; the last ratio is $q=23/80$.

The ridge comparison carries over to the signed factorization.
The calculation of Lemma~\ref{sd:ridge-comparison}, with
$\Sigma H_I\Sigma$ and $\Sigma b_i$, gives
\begin{align*}
 m_i^T(M^TM+\kappa I)^{-1}m_i
 &\leq(\Sigma b_i)^T
    (\Sigma H_I\Sigma+\kappa\Lambda^{-1})^{-1}(\Sigma b_i)\\
 &=b_i^T(H_I+\kappa\Lambda^{-1})^{-1}b_i.
\end{align*}
Here $\Sigma$ commutes with $\Lambda$; the same limiting argument
handles a singular $H_J$.
Set $\kappa=q^2\theta^2/k$ and $\rho=\kappa/\theta=q^2\theta/k$,
and additionally require
\[
 (2/\rho)\|N_{sb}\Lambda_b^{-1/2}\|_F^2\leq1/8.
\]
On these events, $\kappa\leq s_k(M)^2/k$, and the block comparison
in Proposition~\ref{sd:psd-partition} gives
$H_I+\kappa\Lambda^{-1}\geq\Lambda_b/4\oplus(\rho/2)I_s$.
Thus
\[
 \alpha_i\leq8\ell_i+(4/\rho)|b_i^s|^2
       \leq\frac{4k}{q^2\theta}(|A|)_{ii}.
\]

All second moments used in Proposition~\ref{sd:psd-partition}
depend only on $|b_i|^2\leq d$ and are unchanged.
Markov's inequality bounds the failure probabilities of the
three events by
$64kd/\theta$, $16kd/\theta$, and $4k^2d/(q^2\theta)$.
Their sum is at most
\[
 \frac{80k+(25600/529)k^2}{\gamma'_k}<1,\qquad k=1,2,3.
\]
The three ratios are less than $0.988$, $0.983$, and $0.994$.
Also $4k/(q^2\gamma'_k)<1/2$ in every case.
A good partition therefore exists.  Since $Z_k\geq(q\theta)^{2k}$,
\eqref{sd:volume-decay}--\eqref{sd:volume-density}
prove the theorem.
\end{proof}

\begin{remark}\label{sd:scope}
Condition \eqref{sd:absolute-spread} concerns the diagonal of
$|A|$, not merely that of $A$.  The positive and negative spectral
contributions to $a_{ii}$ can cancel for an indefinite matrix.
Such cancellation is absent from $(|A|)_{ii}$.
Thus Theorem~\ref{sd:indefinite} does not settle the unrestricted
indefinite case.
\end{remark}

\subsection{A smaller constant for a separated nonzero spectrum}
The following argument gives a smaller constant when every nonzero
eigenvalue, rather than just the third, is bounded away from zero.

\begin{lemma}[A spread spectral frame]\label{sd:frame}
Let $R$ be an $n\times r$ matrix with $R^TR=I_r$ and let
$\Lambda$ be a nonsingular real diagonal $r\times r$ matrix.
If every row $r_i^T$ of $R$ satisfies $|r_i|^2\leq1/2$, then
for $Q=X^TR\Lambda R^TX+b^TX+c$,
\begin{equation}\label{sd:frame-fourier}
 |\varphi_Q(t)|
 \leq(4\pi)^{r/2}|t|^{-r/2}|\det\Lambda|^{-1/2}.
\end{equation}
\end{lemma}

\begin{proof}
Put $c_i=|r_i|^2$ and $u_i=r_i/\sqrt{c_i}$ when $c_i>0$.
Then $\sum_i c_i u_i u_i^T=I_r$ and $\sum_i c_i=r$.
The geometric Brascamp--Lieb inequality \citep{BL1976,Lieb1990}
applied to
$F_i(s)=|\varphi_{X_i}(\sqrt{c_i}s+tb_i)|^{1/c_i}$ gives
\begin{align*}
 \int_{\mathbb R^r}
  \prod_i|\varphi_{X_i}(\langle r_i,u\rangle+tb_i)|\,du
 &\leq\prod_{i:c_i>0}\left(\int_{\mathbb R}F_i(s)\,ds\right)^{c_i}\\
 &\leq(4\pi)^r.
\end{align*}
The last inequality is Lemma~\ref{four:lp} with exponent
$1/c_i\geq2$, after changing variables.  Rows with $c_i=0$
contribute factors of modulus at most one.

For a justified quadratic Fourier inversion, put $Y=R^TX$ and
regularize the quadratic exponential as
\[
 e^{itY^T\Lambda Y-\varepsilon|Y|^2},\qquad \varepsilon>0.
\]
The Fourier transform with kernel $e^{-iu^Ty}$ of this complex
Gaussian has modulus at most
$\pi^{r/2}|t|^{-r/2}|\det\Lambda|^{-1/2}$.
Fourier inversion and independence therefore bound the
regularized expectation by
\[
 (4\pi|t|)^{-r/2}|\det\Lambda|^{-1/2}
 \int_{\mathbb R^r}
   \prod_i|\varphi_{X_i}(\langle r_i,u\rangle+tb_i)|\,du.
\]
For $\varepsilon>0$ the Gaussian transform is integrable, which
justifies the initial use of Fubini.  The preceding
Brascamp--Lieb bound supplies an integrable majorant for the
product after multiplication by the uniformly bounded Gaussian
Fourier factor.  Let $\varepsilon\downarrow0$ to prove
\eqref{sd:frame-fourier}.
\end{proof}

\begin{theorem}[Separated nonzero semidefinite spectrum]\label{sd:separated}
Suppose $A$ is positive semidefinite of rank $r\geq3$ and every
nonzero eigenvalue is at least $\theta>0$.  Then
\begin{equation}\label{sd:separated-bound}
 \|p_Q\|_\infty\leq130/\theta
\end{equation}
for all $b,c$.  The same assertion holds if all nonzero
eigenvalues are at most $-\theta$.
\end{theorem}

\begin{proof}
Scale so that $\theta=1$, and write $A=R\Lambda R^T$ with
$R^TR=I_r$ and $\Lambda\geq I_r$.  Then $|r_i|^2\leq a_{ii}$.
Use a positive pivot if some $a_{ii}\geq1/2$.
Otherwise Lemma~\ref{sd:frame} applies, and
\[
 |\varphi_Q(t)|
 \leq\min(1,(4\pi/|t|)^{r/2})
 \leq\min(1,(4\pi/|t|)^{j/2})\quad(r\geq j).
\]
For rank at least one, a pivot gives $\alpha_1(Q)\leq4\sqrt2<15$.
Without a pivot, Lemma~\ref{sd:fejer} gives
$\alpha_1(Q)\leq2\sqrt2\sqrt{4\pi}/\sin^2 1<15$.
Thus $\alpha_1\leq15$ throughout this spectral class.

For rank at least two, a pivot leaves a Schur complement of rank
one less with all nonzero eigenvalues at least one.
The first level and Lemma~\ref{sd:composition}(2) give
$\alpha_2(Q)\leq4\sqrt2\sqrt2\cdot15=120$.
Without a pivot, Lemma~\ref{sd:fejer} gives
$\alpha_2(Q)\leq4\pi/\sin^2 1<18$.

For rank at least three, a pivot and the second level give
\[
 \|p_Q\|_\infty\leq3\sqrt2\sqrt{120(3+\log120)}<130.
\]
Without a pivot, integrating
$\min(1,(4\pi/|t|)^{3/2})$ gives $\|p_Q\|_\infty\leq12$.
Scaling and reflection complete the proof.
\end{proof}

\subsection{Sharper bounds in definite rank three}

\begin{theorem}[Exact definite rank three]\label{sd:rank-three}
Suppose $A$ has rank three and its nonzero eigenvalues are all
at least $\theta>0$, or all at most $-\theta$.  Then, for every
$b\in\mathbb R^n$ and $c\in\mathbb R$,
\begin{equation}\label{sd:rank-three-bound}
 \|p_{X^TAX+b^TX+c}\|_\infty\leq58/\theta.
\end{equation}
If in addition $b\in\operatorname{range}(A)$, the stronger bound
\begin{equation}\label{sd:rank-three-range}
 \|p_{X^TAX+b^TX+c}\|_\infty\leq3e/\theta
\end{equation}
holds.
\end{theorem}

\begin{proof}
Reflection and scaling reduce the assertion to $A\geq0$,
$\operatorname{rank}(A)=3$, and $\lambda_3(A)\geq1$.
Write $A=R\Lambda R^T$ with $R^TR=I_3$ and $\Lambda\geq I_3$.
If every row $r_i^T$ of $R$ satisfies $|r_i|^2\leq1/2$,
Theorem~\ref{four:diffuse}, with the individual linear terms as
separable phases, gives
\[
 \|p_Q\|_\infty\leq3e\,(\det\Lambda)^{-1/3}\leq3e.
 \tag*{}
\]
This proves both assertions in this case.

Otherwise choose $i$ with $|r_i|^2>1/2$.  Then
$a=a_{ii}\geq|r_i|^2>1/2$.
Lemma~\ref{sd:pivot} writes
$Q=a(X_i-m(X_{-i}))^2+T(X_{-i})$.
The matrix of $T$ has rank two, with both nonzero eigenvalues
at least one; its linear term may be arbitrary.
Theorem~\ref{thm:rank-two-affine} therefore gives
$\alpha_2(T)\leq46$.  The conditional square density yields
\[
 \|p_Q\|_\infty
 \leq \sqrt2\sup_s\mathbb E|s-T|^{-1/2}.
 \tag*{}
\]
For any $V\geq1$, the layer-cake calculation used in
Lemma~\ref{sd:composition} gives
\[
 \mathbb E|s-T|^{-1/2}
 \leq V+\frac{46(3+2\log V)}{V}.
 \tag*{}
\]
Taking $V=20$ proves \eqref{sd:rank-three-bound}, since
\[
 \sqrt2\left(20+\frac{46(3+2\log20)}{20}\right)
 <57.531<58.
 \tag*{}
\]

Now suppose $b\in\operatorname{range}(A)$.  Choose $z$ with
$2Az=b$ and replace $X$ by $X+z$, absorbing the resulting
constant into $c$.  Translation preserves independence and all
coordinate density caps, so it is enough to treat $b=0$.
In the heavy-row case the Schur complement is a positive
semidefinite matrix of rank two with nonzero eigenvalues
$\mu_1,\mu_2\geq1$.  Thus
\[
 T=Y^T\operatorname{diag}(\mu_1,\mu_2)Y+c',
 \qquad Y=U^TX_{-i},\quad U^TU=I_2.
 \tag*{}
\]
Lemma~\ref{four:projection} bounds the joint density of $Y$ by
$2$.  Lemma~\ref{geom:definite} consequently gives
$\|p_T\|_\infty\leq2\pi$.
Keep the support restriction in the square density formula:
\[
 p_Q(s)\leq a^{-1/2}
      \mathbb E\bigl[(s-T)^{-1/2}{\bf1}_{\{T<s\}}\bigr]
 \quad\hbox{for almost every }s.
 \tag*{}
\]
For any real random variable with density bounded by $M$,
layer cake gives
\[
 \mathbb E\bigl[(s-T)^{-1/2}{\bf1}_{\{T<s\}}\bigr]
 =\int_0^\infty\mathbb P(s-t^{-2}<T<s)\,dt
 \leq\int_0^\infty\min(1,Mt^{-2})\,dt=2\sqrt M.
 \tag*{}
\]
It follows that $\|p_Q\|_\infty\leq4\sqrt\pi<3e$.
Together with the diffuse case, this proves
\eqref{sd:rank-three-range}.  Restoring the scale $\theta$
completes the proof.
\end{proof}
